\chapter{Աշխատանքի մեկնարկ}
\label{ch:start}

\section{Առաջին հանդիպման ստուգաթերթ}
\begin{itemize}
  \item Ընդհանուր git ռեպոզիտորիա և համատեղ թղթապանակ։
  \item AUP ZU3 հարթակի հանձնում (ուսանող 1), լաբորատորիայի մուտք և անվտանգության հրահանգավորում։
  \item Սերվերային համակարգչի կամ վիրտուալ մեքենայի հասանելիություն (ուսանող 2)։
  \item Հաղորդակցման ալիք և շաբաթական հանդիպման ժամ։
\end{itemize}

\section{Միջավայրի պատրաստում}
\begin{enumerate}
  \item \textbf{Ուսանող 1.} տեղադրել Vivado և Vitis (նույն տարբերակը, որ նշված է ռեպոզիտորիայում), ավելացնել AUP ZU3-ի board files-ը, կառուցել և բեռնել «hello world» նախագիծ։
  \item \textbf{Ուսանող 2.} տեղադրել Python, PostgreSQL և ընտրված frontend գործիքները, գործարկել դատարկ API և dashboard։
\end{enumerate}

\section{Ներածական վարժություններ}
\begin{warmup}{Counter և testbench (ուսանող 1)}{6 ժ}
  Նախագծել կոնֆիգուրելի counter, որը պարամետրով ընտրում է քայլը, և գրել ինքնաստուգվող testbench։
  \donewhen{սիմուլյացիան բոլոր դեպքերում PASS է ցույց տալիս, կոդը ռեպոզիտորիայում է։}
\end{warmup}

\begin{warmup}{LED՝ AXI4-Lite ռեգիստրով (ուսանող 1)}{10 ժ}
  Ստեղծել AXI4-Lite IP մեկ ռեգիստրով, որը կառավարում է LED-ի թարթման հաճախությունը, և փոխել այն PS-ի ծրագրից։
  \donewhen{ռեգիստրում գրված արժեքը փոխում է թարթման հաճախությունը հարթակի վրա։}
\end{warmup}

\begin{warmup}{Տվյալների ընդունում և գրաֆիկ (ուսանող 2)}{8 ժ}
  Գրել API endpoint, որը ընդունում է բաժին~\ref{ch:data}-ի JSON-ը, պահում է բազայում, և սկրիպտ, որը վայրկյանում 10 սինուսոիդային նմուշ է ուղարկում։
  \donewhen{dashboard-ը իրական ժամանակում ցույց է տալիս ազդանշանը, տվյալները պահված են բազայում։}
\end{warmup}

\begin{warmup}{Միջերեսի համաձայնեցում (երկուսն էլ)}{4 ժ}
  Համատեղ լրացնել Աղյուսակ~\ref{tab:ifaces}-ը և գրել JSON հաղորդագրությունների օրինակներ ռեպոզիտորիայում։
  \donewhen{ղեկավարները հաստատել են միջերեսի փաստաթուղթը։}
\end{warmup}

\section{Առաջին ընթերցանություն}
\textbf{Ուսանող 1.} \cite{zynqbook} (PS/PL և AXI գլուխները), ապա \cite{axispec}-ի AXI4-Lite մասը։
\textbf{Ուսանող 2.} \cite{fastapi} (Tutorial), \cite{kleppmann2017}, գլ.~2։
